\documentclass[addpoints,10pt]{exam}

\usepackage{amsmath,amsthm,enumitem,wrapfig,amsfonts,mathtools}
\usepackage[mathscr]{euscript}
\usepackage[super]{nth}
\usepackage{dsfont}
\usepackage{xparse}
\usepackage{cancel}

\usepackage{geometry}
\usepackage[T1]{fontenc}
\renewcommand{\rmdefault}{ptm}
\usepackage{amsmath,amsfonts,amsthm,amssymb}
\usepackage{bm}
\usepackage{graphicx}
\usepackage{sectsty}
\usepackage{pgfplots}
\pgfplotsset{compat=1.18}
\usetikzlibrary{arrows.meta}
\usepackage{xcolor}
\definecolor{darkpastelgreen}{rgb}{0.01, 0.75, 0.24}
\definecolor{blue-violet}{rgb}{0.54, 0.17, 0.89}

\usepackage{polynom}
\newcounter{cprob}
\newenvironment{cprob}[1]{%
    \setcounter{cprob}{#1}%
    \noindent\textbf{Problem \thecprob.}%
}{%
    \par\bigskip%
}

\theoremstyle{plain}
\newcommand{\theoremname}{Theorem}
\newtheorem{thm}{\protect\theoremname}
\theoremstyle{definition}
\newtheorem{prob}[thm]{Problem}
\newtheorem*{problem*}{Open Problem}
\theoremstyle{plain}
\newtheorem{conjecture}[thm]{Conjecture}
\theoremstyle{plain}
\newtheorem{lem}[thm]{Lemma}
\newtheorem*{lem*}{Lemma}
\newtheorem{obs}[thm]{Observation}
\newtheorem{cor}[thm]{Corollary}
\theoremstyle{definition}
\newtheorem{definition}[thm]{Definition}

\let\oldprob\prob
\let\endoldprob\endprob
\renewenvironment{prob}
  {\begin{singlespace}\oldprob}
  {\endoldprob\end{singlespace}}

\newcommand{\belowtitle}{\leavevmode\newline}
\newcommand{\Observe}{\text{Observe.}}
\newcommand{\IF}{\mathbf{(\Rightarrow)}}
\newcommand{\FI}{\mathbf{(\Leftarrow)}}
\newcommand{\class}[2][]{\ensuremath{\left[\,#2\,\right]_{#1}}}

\newcommand{\horrule}[1]{\rule{\linewidth}{#1}}
\newcommand{\kkk}{\ensuremath{\Bbbk}}
\newcommand{\CC}{\ensuremath{\mathbb{C}}}
\newcommand{\FF}{\ensuremath{\mathbb{F}}}
\newcommand{\KK}{\ensuremath{\mathbb{K}}}
\newcommand{\NN}{\ensuremath{\mathbb{N}}}
\newcommand{\QQ}{\ensuremath{\mathbb{Q}}}
\newcommand{\RR}{\ensuremath{\mathbb{R}}}
\newcommand{\ZZ}{\ensuremath{\mathbb{Z}}}
\newcommand{\MM}{\ensuremath{\mathcal{M}}}
\newcommand{\TT}{\ensuremath{\mathcal{T}}}
\newcommand{\BB}{\ensuremath{\mathcal{B}}}
\newcommand{\VV}{\ensuremath{\mathcal{V}}}
\newcommand{\WW}{\ensuremath{\mathcal{W}}}
\newcommand{\UU}{\ensuremath{\mathcal{U}}}
\newcommand{\PP}{\ensuremath{\mathcal{P}}}
\newcommand{\LL}{\ensuremath{\mathcal{L}}}
\newcommand{\kk}{\ensuremath{\mathds{k}}}
\newcommand{\EE}{\ensuremath{\mathbb{E}}}
\newcommand{\Ss}{\ensuremath{\mathcal{S}}}

\newcommand{\sm}{\char`\\}

\DeclarePairedDelimiter{\ip}{\langle}{\rangle}
\DeclarePairedDelimiter{\norm}{\lVert}{\rVert}
\DeclarePairedDelimiter{\sqb}{\lbrack}{\rbrack}

\newcommand{\floor}[1]{\left\lfloor #1 \right\rfloor}
\newcommand{\ceil}[1]{\left\lceil #1 \right\rceil}
\newcommand{\mbf}[1]{\ensuremath{\mathbf{#1}}}
\newcommand{\tbf}[1]{\textbf{ #1 }}
\newcommand{\Span}{\ensuremath{\mathrm{Span}}}
\newcommand{\Char}[1]{\mathrm{Char}\; #1}
\DeclareMathOperator{\lcm}{lcm}
\newcommand{\id}{\ensuremath{\mathrm{id}}}
\newcommand{\Gal}[2]{\mathrm{Gal}(#1/#2)}
\newcommand{\Aut}{\ensuremath{\mathrm{Aut}}}
\newcommand{\Fix}[2]{\mathrm{Fix}_{#1}(#2)}
\newcommand{\nil}{\ensuremath{\mathrm{nil}}}
\newcommand{\Hom}{\ensuremath{\mathrm{Hom}}}
\newcommand{\Obj}{\ensuremath{\mathrm{Obj}}}

\newcommand{\Rng}{\ensuremath{\mathrm{Rng}}}
\newcommand{\rank}{\ensuremath{\operatorname{rank}}}
\newcommand{\Tor}{\ensuremath{\operatorname{Tor}}}

\makeatletter
\renewcommand*\env@matrix[1][*\c@MaxMatrixCols c]{%
  \hskip -\arraycolsep
  \let\@ifnextchar\new@ifnextchar
  \array{#1}}
\makeatother

\def\env@matrix{\hskip -\arraycolsep
  \let\@ifnextchar\new@ifnextchar
  \array{*\c@MaxMatrixCols c}}

\newcommand{\proj}[2]{\text{proj}_{#1}(#2)}

%%% Formatting: Page Header
\newcommand{\StudentName}{Danny Banegas}
\newcommand{\AssignmentName}{Homework 13}
\newcommand{\CourseName}{MATH 720 - Algebra I}

\pagestyle{headandfoot}
\runningheadrule
\firstpageheadrule
\firstpageheader{\CourseName}{\StudentName}{\AssignmentName}
\runningheader{\CourseName}{\StudentName}{\AssignmentName}
\firstpagefooter{}{\thepage}{}
\runningfooter{}{\thepage}{}

\printanswers

\DeclareMathAlphabet{\mathcal}{OMS}{cmsy}{m}{n}

\usepackage{parskip}
\usepackage{setspace}

\begin{document}

%%%%%%%%%%%%%%%%%%%%%%%%%%%%%%%%%%%%%%%%%%%%%%%%% 1 %%%%%%%%%%%%%%%%%%%%%%%%%%%%%%%%%%%%%%%%
\setcounter{thm}{0}

\begin{prob}
Consider the matrix
\[
\begin{pmatrix}
9 & 0\\
0 & 15
\end{pmatrix}.
\]

Write a sequence of $\ZZ$-linear row and column operations to get the matrix
\[
\begin{pmatrix}
3 & 0\\
0 & 45
\end{pmatrix}.
\]

Then explain a general method to transform a matrix
\[
\begin{pmatrix}
a & 0\\
0 & b
\end{pmatrix}
\]
to
\[
\begin{pmatrix}
\gcd(a,b) & 0\\
0 & \lcm(a,b)
\end{pmatrix},
\text{ whenever $a,b$ are elements of a PID $R$.}\]

\end{prob}

\begin{proof}
First, observe that $\gcd(a,b)=\frac{ab}{\lcm(a,b)}$. Let
$d=\gcd(a,b)$. Then $a=\alpha d$ and $b=\beta d$ for some
$\alpha,\beta\in R$. Also, $\lcm(a,b)=\alpha\beta d$.

Now suppose $\delta\mid \alpha,\beta$. Then
$a=\alpha d$ and $b=\beta d$, so $\delta d\mid a,b$. Hence
$\delta d\mid \gcd(a,b)=d$, so $\delta$ is a unit. Thus
$\gcd(\alpha,\beta)=1$. Since $R$ is a PID, there exist $x,y\in R$
such that $x\alpha+y\beta=1$. Multiplying by $d$, we get
$xa+yb=d$.

We use this relation to transform the matrix. Starting with
$\begin{pmatrix} a&0\\0&b\end{pmatrix}$, add column $2$ to column $1$
to get $\begin{pmatrix} a&0\\b&b\end{pmatrix}$. Then use the Euclidean
process on the entries in column $1$ until the first column becomes
$\begin{pmatrix} d\\0\end{pmatrix}$. Then we get a matrix of the form
$\begin{pmatrix} d&u\\0&v\end{pmatrix}$. Then subtract multiples of
column $1$ from column $2$ to clear $u$, obtaining
$\begin{pmatrix} d&0\\0&v\end{pmatrix}$. We can do this since $u$ was obtained from multiples of $a$ and $b$ both of whom share $d$ as a divisor.

Every row and column operation preserves the ideal generated by all
entries, so $d=\gcd(a,b)$ remains the first diagonal entry. Also row and
column operations change determinant only by multiplication by a unit.
Therefore, $ab$ is associate to $dv$ and $v$ is associate to
$\frac{ab}{d}=\lcm(a,b)$. So the final matrix is equivalent to
$\begin{pmatrix} \gcd(a,b)&0\\0&\lcm(a,b)\end{pmatrix}$.

For the specific matrix,
\begin{align*}
&\begin{pmatrix}
9&0\\
0&15
\end{pmatrix}
\overset{C_2= C_1+C_2}{\longrightarrow}
\begin{pmatrix}
9&0\\
15&15
\end{pmatrix}
\overset{R_2= R_2-R_1}{\longrightarrow}
\begin{pmatrix}
9&0\\
6&15
\end{pmatrix}
\overset{R_1= R_1-R_2}{\longrightarrow}
\begin{pmatrix}
3&-15\\
6&15
\end{pmatrix}
\overset{R_2= R_2-R_1}{\longrightarrow}
\begin{pmatrix}
3&-15\\
3&30
\end{pmatrix}\\
&\overset{R_1= R_1-R_2}{\longrightarrow}
\begin{pmatrix}
0&-45\\
3&30
\end{pmatrix}
\overset{R_1\leftrightarrow R_2}{\longrightarrow}
\begin{pmatrix}
3&30\\
0&-45
\end{pmatrix}
\overset{C_2=C_2-10C_1}{\longrightarrow}
\begin{pmatrix}
3&0\\
0&-45
\end{pmatrix}
\overset{R_2=-R_2}{\longrightarrow}
\begin{pmatrix}
3&0\\
0&45
\end{pmatrix}.
\end{align*}
\end{proof}

%%%%%%%%%%%%%%%%%%%%%%%%%%%%%%%%%%%%%%%%%%%%%%%%% 2 %%%%%%%%%%%%%%%%%%%%%%%%%%%%%%%%%%%%%%%%

\begin{prob}
Read the full proof of Smith Normal Form (in Smith Normal Form Supplement under Week 16), and then list any steps which require $R$ to be a PID.
\end{prob}

The proof uses that $R$ is a PID in the following places.

First, immediately after defining the function that counts the number of
prime factors $\ell(a)$ of an element $a\in R$, this only makes sense in
full generality if every element has a factorization into primes. This is
true because every PID is a UFD.

Second, the proof uses that every ideal generated by finitely many
elements is principal. For example, in Case 2a it says that since $R$ is
a PID, $\langle a_1,a_2\rangle=\langle d\rangle$. This step requires
$R$ to be a PID.

Third, the same issue occurs later whenever the proof replaces an ideal
generated by several entries with a single generator.
%%%%%%%%%%%%%%%%%%%%%%%%%%%%%%%%%%%%%%%%%%%%%%%%% 3 %%%%%%%%%%%%%%%%%%%%%%%%%%%%%%%%%%%%%%%%

\newpage

\begin{prob}
Let
\[
R=k[x,y,z].
\]

Explain what goes wrong when you try to find a ``Smith Normal Form'' of
\[
\begin{pmatrix}
x & y\\
z & x+y
\end{pmatrix}.
\]
\end{prob}

\begin{proof}
Let $R=k[x,y,z]$. Notice that $R$ is not a PID. In particular,
$\langle x,y,z\rangle$ is not principal since if
$\langle x,y,z\rangle=\langle d\rangle$, then $d$ divides $x$, $y$, and
$z$. Since $x,y,z$ have no common nonunit factor, $d$ must be a unit.
So $\langle d\rangle=R$. But this would imply
$\langle x,y,z\rangle=R$, which is impossible since every element of
$\langle x,y,z\rangle$ has no constant term.

Now consider
$\begin{pmatrix}x&y\\ z&x+y\end{pmatrix}$. The entries generate the
ideal $\langle x,y,z,x+y\rangle=\langle x,y,z\rangle$. If a Smith normal
form existed by the usual PID algorithm, the first diagonal entry would
generate the ideal generated by all entries. But
$\langle x,y,z\rangle$ is not principal. Therefore there is no single
gcd of the entries obtained as an $R$-linear combination of
$x,y,z,x+y$.

So the first step of the Smith normal form algorithm already fails:
we cannot replace the entries by a greatest common divisor using row and
column operations. Therefore the usual Smith normal form process does not
work over $k[x,y,z]$.

\end{proof}
%%%%%%%%%%%%%%%%%%%%%%%%%%%%%%%%%%%%%%%%%%%%%%%%% 4 %%%%%%%%%%%%%%%%%%%%%%%%%%%%%%%%%%%%%%%%

\newpage

\begin{prob}
Let $R$ be an integral domain with field of fractions $F=Q(R)$, and let
$M$ be any $R$-module.

Prove that the rank of $M$ is the dimension of the vector space
$F\otimes_R M$ over $F$.
\end{prob}

\begin{proof}
Let $F=Q(R)$. $(\implies)$ Suppose $m_1,\ldots,m_n\in M$ are $R$-linearly independent.
We prove that $1\otimes m_1,\ldots,1\otimes m_n$ are $F$-linearly
independent in $F\otimes_R M$.

Suppose $\sum_{i=1}^n \lambda_i(1\otimes m_i)=0$, where
$\lambda_i\in F$. Write each $\lambda_i=a_i/s_i$ with $a_i\in R$ and
$s_i\in R\setminus\{0\}$. Let $s=s_1\cdots s_n$. Multiplying by $s$,
we get $\sum_{i=1}^n b_i(1\otimes m_i)=0$ for some $b_i\in R$.
Equivalently, $1\otimes \sum_{i=1}^n b_i m_i=0$. Since the $m_i$ are
$R$-linearly independent, each $b_i=0$, and therefore each
$\lambda_i=0$. So the pure tensors are $F$-linearly independent.

$(\impliedby)$ On the other hand, suppose $1\otimes m_1,\ldots,1\otimes m_n$ are
$F$-linearly independent. If $\sum_{i=1}^n r_i m_i=0$ with $r_i\in R$,
then tensoring gives $\sum_{i=1}^n r_i(1\otimes m_i)=0$ in
$F\otimes_R M$. Since the tensors are $F$-linearly independent, each
$r_i=0$. This holds since $R_{i}\cong_{R}\frac{R_{i}}{1}\leq_{R} F$. Therefore, $m_1,\ldots,m_n$ are $R$-linearly independent.

So then $R$-linearly independent subsets of $M$ correspond to
$F$-linearly independent subsets of $F\otimes_R M$. That is, if $\rank M\leq n$, there exists some $R$-linearly independent $m_{1},\hdots, m_{n}\in M$. If we suppose $\rank F<n$ we then get a contradiction since $1\otimes m_{1},\hdots, 1\otimes m_{n}$ are $F$-linearly independent. If $M$ has infinite rank, then we can form an infinite set of $R$-linearly independent elements in $M$ and a set of $F$-linearly independent pure tensors corresponding to them. So $\rank M=\dim_{F}(F\otimes_{R} M).$
\end{proof}

%%%%%%%%%%%%%%%%%%%%%%%%%%%%%%%%%%%%%%%%%%%%%%%%% 5 %%%%%%%%%%%%%%%%%%%%%%%%%%%%%%%%%%%%%%%%

\newpage

\begin{prob}
Let $R=\ZZ[x]$ and let $M=(2,x)$.

\begin{enumerate}[label=(\alph*)]
\item Show that $\{2,x\}$ is not a basis for $M$.

\item Show that $M$ has rank $1$ but is not free of rank $1$.

\item Write a presentation for $M$.
\end{enumerate}
\end{prob}

\begin{proof}
(a)
In $M=\langle 2,x\rangle\subseteq \ZZ[x]$, we have the nontrivial relation
$x\cdot 2+(-2)\cdot x=0$. The coefficients $x$ and $-2$ are both non-zero
zero in $\ZZ[x]$. So by our equivalence theorem $\langle 2,x\rangle$ cannot be a basis for $\ZZ[x]$

(b)
Since $M$ is a nonzero submodule of $R=\ZZ[x]$, and $R$ is an integral
domain, any two nonzero elements of $M$ are $R$-linearly dependent since for any $m,n\in M\setminus\{0\}$ and any we have that $nm-mn=0$, a non-trivial $R$-linear combination. So no
two elements of $M$ can be linearly independent, and $\rank M\leq 1$ but $M$ does contain a non-zero element $2$ which is then linearly independent by itself and so $\rank M=1$

Now suppose $M$ were free of rank $1$. Then $M=\langle f\rangle$ for
some $f\in M$, so $M$ would be a principal ideal. Since $2,x\in M$, we
would have $f\mid 2$ and $f\mid x$ in $\ZZ[x]$. The only common divisors
of $2$ and $x$ are units, so $f$ would be a unit. Then
$M=\langle f\rangle=R$, which is false because $1\notin (2,x)$. So
$M$ is not free of rank $1$.

(c)
Let $\phi:R^2\to M$ be defined by $\phi(a,b)=2a+xb$. This map is obviously surjective.
$2a+xb=0\implies xb=-2a$. So $2a=-xb$, and $x\mid 2a$. Since $x$ is prime in $\ZZ[x]$ and
$x\nmid 2$, we get that $x\mid a$ and $a=xh$ for some $h\in R$.

Substituting back into $2a=-xb$ gives $2xh=-xb$. Since $R=\ZZ[x]$ is an
integral domain, we can cancel $x$ to get $2h=-b$. So $b=-2h$.

Therefore $(a,b)=(xh,-2h)=h(x,-2)$, so
$\ker\phi\subseteq \langle (x,-2)\rangle$.
Then in the other direction we have that
$\phi(x,-2)=2x+x(-2)=0$. Therefore, $\ker\phi=\langle (x,-2)\rangle$.

Therefore a presentation is
$R\xrightarrow{\begin{pmatrix}x\\-2\end{pmatrix}}R^2\to M\to 0$.
Equivalently, $M\cong R^2/\langle (x,-2)\rangle$.

\end{proof}

%%%%%%%%%%%%%%%%%%%%%%%%%%%%%%%%%%%%%%%%%%%%%%%%% 6 %%%%%%%%%%%%%%%%%%%%%%%%%%%%%%%%%%%%%%%%

\newpage

\begin{prob}
Let $R$ be a ring and let $M$ be an $R$-module.

\begin{enumerate}[label=(\alph*)]
\item Prove that if $M$ is projective, then $\Tor(M)=(0)$.

\item Prove that a finitely generated module over a PID is projective if and only if it is free.
\end{enumerate}
\end{prob}

\begin{proof}
(a)
Let $M$ be projective. Then there exists an $R$-module $N$ such that
$M\oplus N$ is free. Since $M\oplus N$ embeds into a free module, if
$m\in M$ is torsion, then there exists nonzero $r\in R$ such that
$rm=0$. But then $r(m,0)=(0,0)$ in the free module $M\oplus N$. Since
free modules over a domain are torsion-free, $(m,0)=0$, and so $m=0$.
Thus, $\Tor(M)=\langle 0\rangle$.

(b)
If $M$ is free, then $M$ is projective automatically by our equivalence theorem since it is a direct summand of its free self.

On the other hand, suppose $M$ is finitely generated and projective over a PID.
Since $M$ is finitely generated over a PID, the structure theorem gives
$M\cong R^r\oplus R/(d_1)\oplus\cdots\oplus R/(d_k)$, where each
$d_i$ is a nonzero nonunit. By part (a), projective modules are
torsion-free, so the torsion summands must vanish. Therefore
$M\cong R^r$ and $M$ is free.

\end{proof}

%%%%%%%%%%%%%%%%%%%%%%%%%%%%%%%%%%%%%%%%%%%%%%%%% 7 %%%%%%%%%%%%%%%%%%%%%%%%%%%%%%%%%%%%%%%%

\newpage

\begin{prob}
Let $R$ be a PID which is not a field.

Prove that no finitely generated $R$-module is injective.

\textbf{Hint:} Recall that $Q\cong Q_1\oplus Q_2$ is injective if and only if $Q_1$ and $Q_2$ are injective.
\end{prob}

\begin{proof}
We prove that no nonzero finitely generated $R$-module is injective.

Since $R$ is a PID and not a field, choose a nonzero nonunit $r\in R$. First, $R$ is not injective as an $R$-module, or we contradict a previous homework problem where we proved that if $R$ is injective over itself then it must be a field.

Next, let $d\in R$ be a nonzero nonunit. We show $R/\langle d\rangle
$ is not injective. Consider the inclusion $\langle d\rangle\hookrightarrow R$, and let
$h:\langle d\rangle\to R/\langle d\rangle$ be defined by $h(da)=[a]$. This is well-defined since $R$ is a
domain. If $R/\langle d\rangle$ were injective, then $h$ would extend to some
$f:R\to R/\langle d\rangle$. But then
$h(d\cdot 1)=[1]=f(d)=df(1)=[0]$ in $R/\langle d\rangle$. But then $1-0=1\in \langle d\rangle\implies \langle d\rangle=R$, a contradiction since $d$ is a nonzero nonunit$ \implies \langle d\rangle$ is proper. So $R/\langle d\rangle$ is not
injective.

Now let $M$ be a nonzero finitely generated $R$-module. By the structure
theorem for finitely generated modules over a PID,
$M\cong R^n\oplus R/(d_1)\oplus\cdots\oplus R/(d_k)$,
where each $d_i$ is a nonzero nonunit.

If $M$ were injective, then every direct summand of $M$ would be
injective by a previous homework problem. But each possible nonzero summand is either $R$
or some $R/(d_i)$, and we proved neither type is injective.

Therefore no nonzero finitely generated $R$-module is injective.

\end{proof}

%%%%%%%%%%%%%%%%%%%%%%%%%%%%%%%%%%%%%%%%%%%%%%%%% 8 %%%%%%%%%%%%%%%%%%%%%%%%%%%%%%%%%%%%%%%%

\newpage

\begin{prob}
Let
\[
V=\RR[x]/(x^4+2x^2+3x+1).
\]

\begin{enumerate}[label=(\alph*)]
\item Show that $V$ is an $\RR$-vector space with basis
$\beta=\{1,x,x^2,x^3\}$.

\item Prove that the map $T:V\to V$ given by multiplication by $x$ is an $\RR$-linear transformation.

\item Find the matrix $[T]_{\beta}^{\beta}$.
\end{enumerate}
\end{prob}

\begin{proof}
(a)
Let $V=\RR[x]/(x^4+2x^2+3x+1)$. Since $x^{4}=-2x^2-3x-1$ in $V$, every element of $V$ can be expressed with a cubic representative, since that relation reduces all powers higher than $3$. So then $\{1,x,x^{2},x^{3}\}$ (modulo $x^4+2x^2+3x+1$) generates $V$.

If $a_0+a_1x+a_2x^2+a_3x^3=0$ in $V$, then
$a_0+a_1x+a_2x^2+a_3x^3$ belongs to the ideal
$(x^4+2x^2+3x+1)$. But a nonzero element of that ideal has degree at
least $4$, while this polynomial has degree at most $3$. So then all
$a_i=0$ otherwise a degree $1\leq n\leq 3$ polynomial or a constant besides $0$ belongs to $(x^4+2x^2+3x+1)$, both contradictions. So then by our equivalence theorem, $\beta=\{1,x,x^2,x^3\}$ must be a basis for $V$ over $\RR$, and $V\cong \RR^4$ as vector spaces.

(b)
Let $T:V\to V$ be defined by $T(v)=xv$. $\forall u,v\in V,\forall c\in\RR$, we have
$T(u+v)=x(u+v)=xu+xv=T(u)+T(v)$ and
$T(cv)=x(cv)=c(xv)=cT(v)$. Therefore $T$ is $\RR$-linear.

(c)
We compute the images of the basis elements:
$T(1)=x$, $T(x)=x^2$, $T(x^2)=x^3$, and
$T(x^3)=x^4$. In $V$, we have
$x^4=-(2x^2+3x+1)=-1-3x-2x^2$.

Therefore, 
$[T(1)]_\beta=\begin{pmatrix}0\\1\\0\\0\end{pmatrix}$,
$[T(x)]_\beta=\begin{pmatrix}0\\0\\1\\0\end{pmatrix}$,
$[T(x^2)]_\beta=\begin{pmatrix}0\\0\\0\\1\end{pmatrix}$, and
$[T(x^3)]_\beta=\begin{pmatrix}-1\\-3\\-2\\0\end{pmatrix}$.

and therefore
\[
[T]_\beta^\beta=
\begin{pmatrix}
0&0&0&-1\\
1&0&0&-3\\
0&1&0&-2\\
0&0&1&0
\end{pmatrix}.
\]

\end{proof}

\end{document}